\section{Machine Learning Models}
\label{sec:models}

The task is \textbf{regression, not classification}. Four continuous
school-level outcomes are predicted, each on its own: two grade averages on a
0--100 scale (\texttt{math\_avg\_grade}, \texttt{english\_avg\_grade}) and two
advanced-track (5-unit) participation rates on a 0--1 scale
(\texttt{math\_5unit\_participation}, \texttt{english\_5unit\_participation}).

\textbf{Four model families compete on equal terms.} Rather than commit to one
algorithm, we span the useful complexity range from a regularised linear fit to
modern boosted trees (Table~\ref{tab:models}). Every family is tuned within the
same inner folds using a predefined model-specific search space and up to 25
candidate configurations. Ridge's seven regularisation values are evaluated
exhaustively. An earlier version of this study tuned only the boosted-tree model
and left the others at defaults, which makes any conclusion about a winner an
artefact of unequal effort rather than a property of the data.

\begin{table}[H]
\centering
\footnotesize
\renewcommand{\arraystretch}{1.15}
\caption{The four competing model families. All are tuned under the same nested
protocol using model-specific hyperparameter spaces.}
\label{tab:models}
\begin{tabular}{@{}>{\raggedright\arraybackslash}p{2.2cm}>{\raggedright\arraybackslash}p{1.9cm}>{\raggedright\arraybackslash}p{9.7cm}@{}}
\toprule
\textbf{Model} & \textbf{Family} & \textbf{Why it fits this problem} \\
\midrule
Ridge & Linear (L2) & Stable linear baseline the trees must beat. L2 shrinkage handles the many one-hot dummies well. \\
SGD linear & Linear (SGD) & Second linear fit by stochastic gradient descent, a check that the linear result is not a solver artefact. \\
Random Forest & Bagged trees & Captures non-linearities and interactions without scaling. Bagging resists overfitting. \\
Hist\-Gradient\-Boosting & Boosted trees & Efficient boosted-tree model capturing nonlinear relationships and feature interactions in encoded tabular data. \\
\bottomrule
\end{tabular}
\end{table}

\textbf{Nested cross-validation keeps the evaluation honest.} Two kinds of
leakage matter. Structurally, each school contributes up to four year-rows
(2013--2016), so folds are formed by school using GroupKFold(\texttt{semel}).
More subtly, if feature selection or tuning sees the whole dataset then the
evaluation folds have already influenced the model and scores become optimistic.
We therefore \emph{nest}: an outer GroupKFold supplies untouched evaluation
folds, and within each outer \emph{training} fold only we run VIF pruning, Boruta
selection, and an inner GroupKFold search. Every number in
Table~\ref{tab:leaderboard} is an outer-fold result, reported as a mean and
standard deviation across folds.

\begin{table}[H]
\centering
\footnotesize
\renewcommand{\arraystretch}{1.3}
\caption{Comparative model performance: outer-fold $R^2$ (mean $\pm$ standard
deviation over five folds), each family tuned under the same protocol. Random
Forest leads on every target, though the families sit within roughly one
standard deviation of each other on the grade outcomes.}
\label{tab:leaderboard}
\begin{tabular}{@{}>{\raggedright\arraybackslash}p{4.5cm}cccc@{}}
\toprule
\textbf{Target} & \textbf{Ridge} & \textbf{SGD} & \textbf{Rand.\ Forest} & \textbf{HistGB} \\
\midrule
\texttt{math\_avg\_grade}              & $0.405\!\pm\!0.066$ & $0.402\!\pm\!0.058$ & $\mathbf{0.420\!\pm\!0.050}$ & $0.414\!\pm\!0.054$ \\
\texttt{english\_avg\_grade}           & $0.410\!\pm\!0.045$ & $0.410\!\pm\!0.043$ & $\mathbf{0.452\!\pm\!0.042}$ & $0.441\!\pm\!0.043$ \\
\texttt{math\_5unit\_participation}    & $0.323\!\pm\!0.043$ & $0.323\!\pm\!0.041$ & $\mathbf{0.395\!\pm\!0.037}$ & $0.378\!\pm\!0.030$ \\
\texttt{english\_5unit\_participation} & $0.502\!\pm\!0.036$ & $0.504\!\pm\!0.037$ & $\mathbf{0.557\!\pm\!0.053}$ & $0.539\!\pm\!0.064$ \\
\bottomrule
\end{tabular}
\end{table}

Random Forest is selected for every target, though two observations temper the
ranking. On the grade outcomes all four families lie within about one standard
deviation of one another, so the algorithm matters far less than the features.
The trees separate more clearly on participation, where the outcome is bounded
and the relationships less linear. Feature selection was stable across folds:
11--12 features for the Math targets and 13--15 for English.
